% Part: methods
% Chapter: proofs
% Section: example-1

\documentclass[../../../include/open-logic-section]{subfiles}

\begin{document}

\olfileid{mth}{prf}{ex1}

\olsection{ఒక ఉదాహరణ}

మన మొదటి ఉదాహరణ సమితుల సమ్మేళనాలు, ఛేదనల గురించి ఒక సరళమైన విషయం. ఇది నిర్వచనాలను విప్పడం, సంయోగాలను నిరూపించడం, సార్వత్రిక వాదనలను నిరూపించడం, సందర్భాలవారీ నిరూపణలను వివరిస్తుంది.

\begin{prop}
ఏ సమితులు $A$, $B$, $C$లకైనా $A \cup (B \cap C) = (A \cup B)
\cap (A \cup C)$.
\end{prop}

దాన్ని నిరూపిద్దాం!{}

\begin{proof}
ఏ సమితులు $A$, $B$, $C$లకైనా $A \cup (B \cap
C) = (A \cup B) \cap (A \cup C)$ అని చూపాలి.
\begin{quote}
మొదట ప్రతిపాదనలోని ``$=$'' నిర్వచనాన్ని విప్పుదాం. రెండు సమితులు ఒకటే అని నిరూపించడం అంటే వాటిలో ఒకే \tetoken{మూలకాలు}{!!{element}s} ఉన్నాయని చూపడం. అంటే $A \cup
(B \cap C)$లోని అన్ని \tetoken{మూలకాలు}{!!{element}s}, $(A \cup B) \cap (A \cup C)$లోనూ \tetoken{మూలకాలు}{!!{element}s}; వ్యతిరేక దిశలోనూ అలాగే. ``వ్యతిరేక దిశలోనూ'' అంటే $(A
\cup B) \cap (A \cup C)$లోని ప్రతి \tetoken{మూలకం}{!!{element}}, $A \cup (B \cap
C)$లో \tetoken{ఒక మూలకం}{!!a{element}} కావాలి. నిర్వచనాన్ని విప్పితే ఒక సంయోగాన్ని నిరూపించాల్సి ఉందని తెలుస్తుంది. దీన్ని నమోదు చేద్దాం:
\end{quote}
నిర్వచనం ప్రకారం $A \cup (B \cap C) = (A \cup B) \cap (A \cup C)$ అప్పుడూ అప్పుడే $A \cup (B \cap C)$లోని ప్రతి \tetoken{మూలకం}{!!{element}}, $(A
\cup B) \cap (A \cup C)$లో \tetoken{ఒక మూలకం}{!!a{element}}; అలాగే $(A \cup B) \cap (A
\cup C)$లోని ప్రతి \tetoken{మూలకం}{!!{element}}, $A \cup (B \cap C)$లో \tetoken{ఒక మూలకం}{!!a{element}}.
\begin{quote}
ఇది సంయోగం కనుక ప్రతి సంయోజ్యాన్ని విడిగా నిరూపించాలి. మొదటిదానితో మొదలుపెడదాం: $A \cup (B \cap C)$లోని ప్రతి \tetoken{మూలకం}{!!{element}}, $(A
\cup B) \cap (A \cup C)$లో \tetoken{ఒక మూలకం}{!!a{element}} అని నిరూపిద్దాం.

ఇది సార్వత్రిక వాదన; కాబట్టి $A \cup (B \cap C)$లోని ఏదైనా ఒక \tetoken{మూలకాన్ని}{!!{element}} తీసుకొని, అది $(A \cup B) \cap (A \cup C)$లోనూ \tetoken{ఒక మూలకం}{!!a{element}} అని చూపాలి. ఈ ఏదైనా ఒక \tetoken{మూలకానికి}{!!{element}}~$z$ అనే చరరాశి పేరు పెడదాం. నిరూపణ ఇలా సాగుతుంది:
\end{quote}
మొదట $A \cup (B \cap C)$లోని ప్రతి \tetoken{మూలకం}{!!{element}}, $(A \cup B) \cap (A \cup C)$లోనూ \tetoken{ఒక మూలకం}{!!a{element}} అని నిరూపిద్దాం. $z \in A \cup (B
\cap C)$ అనుకుందాం. $z \in (A \cup B) \cap (A \cup C)$ అని చూపాలి.
\begin{quote}
ఇప్పుడు $\cup$,~$\cap$ నిర్వచనాలను విప్పాలి. ఉదాహరణకు $\cup$ నిర్వచనం: $A \cup B = \Setabs{z}{z \in A
  \text{ or } z \in B}$. ``$A \cup (B
\cap C)$''కు నిర్వచనాన్ని వర్తింపజేస్తే, నిర్వచనంలోని ``$B$'' పాత్రను ఇప్పుడు ``$B \cap C$'' పోషిస్తుంది. కనుక $A \cup (B \cap C) = \Setabs{z}{z \in A \text{ or }
  z \in B \cap C}$. అందుచేత $z \in A \cup (B \cap C)$ అనే పరికల్పన అంటే $z \in \Setabs{z}{z \in A \text{ or } z \in B \cap
  C}$. అలాగే $z \in \Setabs{z}{\dots z\dots}$ అప్పుడూ అప్పుడే \dots $z$ \dots; ఈ సందర్భంలో $z \in A$ లేదా $z \in B \cap C$.
\end{quote}
$\cup$ నిర్వచనం ప్రకారం $z \in A$ లేదా $z \in B \cap C$.
\begin{quote}
ఇది వికల్పం కాబట్టి సందర్భాలవారీ నిరూపణ ఉపయోగపడుతుంది. రెండు సందర్భాలను తీసుకొని, ప్రతి దానిలో మన లక్ష్య నిష్కర్ష, అంటే ``$z \in (A \cup B) \cap (A \cup
C)$'', వస్తుందని చూపాలి.
\end{quote}
సందర్భం 1: $z \in A$ అనుకుందాం.
\begin{quote}
పరికల్పనల నుంచి ఇంకా ఎక్కువ సమాచారం లేదు. కాబట్టి నిష్కర్షలో ఉపయోగించగలిగేది చూద్దాం. $z \in (A \cup B) \cap (A \cup C)$ అని చూపాలి. $\cap$ నిర్వచనం ప్రకారం ఇది చూపాలంటే $z \in (A \cup B) \cap (A \cup C)$లోని రెండు భాగాలలోనూ, అంటే $(A \cup B)$లోనూ $(A \cup C)$లోనూ ఉండాలి. అయితే $z
\in A \cup B$ అప్పుడూ అప్పుడే $z \in A$ లేదా $z \in B$; సందర్భం~1 పరికల్పనగా $z \in A$ ఇప్పటికే ఉంది. అదే విధంగా $B$ స్థానంలో $C$ తీసుకుంటే $z \in A \cup C$. ఈ ఆలోచన లక్ష్యం నుంచి వెనక్కి సాగింది; ఇప్పుడు నిరూపణకు కావలసిన ముందుకు వెళ్లే క్రమంలో రాద్దాం.
\end{quote}
$z \in A$ కనుక $z \in A$ లేదా $z \in B$; అందువల్ల~$\cup$ నిర్వచనం ప్రకారం $z \in A \cup B$. అలాగే $z \in A \cup C$. కనుక~$\cap$ నిర్వచనం ప్రకారం $z \in (A \cup B) \cap (A \cup C)$.
\begin{quote}
దీంతో సందర్భాలవారీ నిరూపణలో మొదటి సందర్భం పూర్తయింది. ఇప్పుడు $z \in B \cap C$ అయ్యే రెండవ సందర్భంలో నిష్కర్షను పొందాలి.
\end{quote}
సందర్భం 2: $z \in B \cap C$ అనుకుందాం.
\begin{quote}
మళ్లీ రెండు సమితుల ఛేదనతో పనిచేస్తున్నాం. $\cap$ నిర్వచనాన్ని వర్తింపజేద్దాం:
\end{quote}
$z \in B \cap C$ కనుక,~$\cap$ నిర్వచనం ప్రకారం $z$ అనేది $B$, $C$ రెండింటిలోనూ \tetoken{ఒక మూలకం}{!!a{element}}.
\begin{quote}
మళ్లీ నిష్కర్షను చూద్దాం. $z$ అనేది $(A \cup B)$, $(A \cup C)$ రెండింటిలోనూ ఉందని చూపాలి. ఇక్కడ కూడా సమాధానం వెంటనే వస్తుంది.
\end{quote}
$z \in B$ కనుక $z \in (A \cup B)$. $z \in C$ కనుక $z \in (A
\cup C)$ కూడా. కాబట్టి $z \in (A \cup B) \cap (A \cup C)$.
\begin{quote}
ఇక్కడ $\cup$, $\cap$ నిర్వచనాలను మళ్లీ వర్తింపజేశాం. వాటిని ఇప్పటికే చెప్పాం; రెండు సమితులలో ఒకదానిలో $z$ ఉంటే వాటి సమ్మేళనంలోనూ ఉంటుందని ఇప్పటికే చూపాం. కాబట్టి ఈసారి అంత వివరంగా చెప్పనవసరం లేదు.

సందర్భాలవారీ నిరూపణలో రెండవ సందర్భమూ పూర్తయింది; ఇప్పుడు మొదటి నిష్కర్షను చెప్పవచ్చు.
\end{quote}
కాబట్టి $z \in A \cup (B \cap C)$ అయితే $z \in (A \cup B) \cap (A \cup C)$.
\begin{quote}
ఇప్పుడు మరో దిశను చూపాలి: $(A \cup B) \cap (A \cup C)$లోని ప్రతి \tetoken{మూలకం}{!!{element}}, $A \cup (B \cap
C)$లో \tetoken{ఒక మూలకం}{!!a{element}}. ముందులాగే, మొదటి సమితిలో ఏదైనా ఒక \tetoken{మూలకాన్ని}{!!{element}} తీసుకొని అది రెండవ సమితిలోనూ ఉంటుందని చూపి ఈ సార్వత్రిక వాదనను నిరూపిస్తాం. ఏమి చేయబోతున్నామో చెప్పుకుందాం.
\end{quote}
ఇప్పుడు $z \in (A \cup B) \cap (A \cup C)$ అనుకుందాం. $z \in A \cup (B \cap C)$ అని చూపాలి.
\begin{quote}
ఇప్పుడు మన పరికల్పన $z \in (A \cup B) \cap (A
\cup C)$. నిరూపణ మొదటి భాగంలో వాడిన అదే~$z$ను ఇక్కడ మళ్లీ వాడటం గందరగోళంగా ఉండకూడదు. ఆ భాగం ముగిసినప్పుడు అక్కడ $z$ గురించి తీసుకున్న పరికల్పనలు ఇక అమల్లో ఉండవు; కాబట్టి $z$ గురించి కొత్త పరికల్పనలు తీసుకోవచ్చు. సందేహమైతే, ఇక్కడి నుంచి వేరే చరరాశి వాడండి.

$\cap$ నిర్వచనం ప్రకారం $z$ అనేది $A \cup B$, $A \cup C$ రెండింటిలోనూ ఉంది. $\cup$ నిర్వచనాన్ని ఇంకా విప్పితే: $z \in A$ లేదా $z \in B$; అలాగే $z \in A$ లేదా $z \in
C$. మళ్లీ సందర్భాలవారీ నిరూపణలా ఉంది---కానీ ``మరియు'' వల్ల గందరగోళం రావచ్చు. ఇక్కడ $z$ అనేది $A$, $B$, $C$లలో ఏదో ఒకదానిలో ఉంటే చాలని అనుకోవచ్చు; అది తప్పు. జాగ్రత్తగా ప్రతి వికల్పాన్ని వరుసగా చూద్దాం.
\end{quote}
$\cap$ నిర్వచనం ప్రకారం $z \in A \cup B$ మరియు $z \in A \cup C$. $\cup$ నిర్వచనం ప్రకారం $z \in A$ లేదా $z \in B$. సందర్భాలను వేరుచూద్దాం.
\begin{quote}
మొదటి వికల్పంపైనే దృష్టి పెడుతున్నాం కనుక $A \cup C$ను విప్పి వచ్చే రెండవ వికల్పాన్ని ఇంకా తీసుకోలేదు. నిజానికి ఇప్పుడది అవసరం లేదు. మొదటి సందర్భం $z \in A$; ఒక సమితిలోని \tetoken{ఒక మూలకం}{!!a{element}}, ఆ సమితిని మరే సమితితో సమ్మేళనం చేసినా \tetoken{ఒక మూలకం}{!!a{element}}. కాబట్టి సందర్భం~1 సులభం:
\end{quote}
సందర్భం 1: $z \in A$ అనుకుందాం. అప్పుడు $z \in A \cup (B \cap
C)$.
\begin{quote}
ఇక రెండవ సందర్భం: $z \in B$. ఇక్కడ రెండవ $\cup$ను విప్పి మరో సందర్భాలవారీ నిరూపణ చేద్దాం:
\end{quote}
సందర్భం 2: $z \in B$ అనుకుందాం. $z \in A \cup C$ కనుక $z \in
A$ లేదా $z \in C$. సందర్భాలను ఇంకా వేరుచూద్దాం:

సందర్భం 2a: $z \in A$. అప్పుడు మళ్లీ $z \in A \cup (B \cap C)$.
\begin{quote}
ఇది కొంచెం వింతగా అనిపించవచ్చు. ఈ సందర్భానికి~$z \in B$ పరికల్పన నిజంగా అవసరం కాలేదు; అది సమస్య కాదు.
\end{quote}
సందర్భం 2b: $z \in C$. అప్పుడు $z \in B$, $z \in C$ రెండూ ఉన్నాయి; కనుక $z \in B \cap C$, తద్వారా $z \in A \cup (B \cap C)$.
\begin{quote}
దీంతో రెండు సందర్భాలవారీ నిరూపణలూ ముగిశాయి; నిరూపణ రెండవ సగం పూర్తయింది.
\end{quote}
కాబట్టి $z \in (A \cup B) \cap (A \cup C)$ అయితే $z \in A \cup (B \cap C)$.
\end{proof}

\end{document}
